\section{From epistemic mechanics to the transfer relation}

The companion paper in this series formalizes the epistemic mechanics of an evidence-governed research substrate \cite{yiu2026epistemic}. Its results fix the ground this paper stands on: claims are indexed by context, evidence, uncertainty and provenance; scientific authority is a partial order that cannot be faithfully collapsed into a scalar; a projection that maps two worlds requiring different canonical updates to the same represented state admits no correct deterministic policy; and---the invariant this paper inherits directly---recorded experience, retrieval priority and representation changes may alter computational access and search cost without conferring evidential authority. Orion, the system both papers describe, persists task episodes, consolidates scoped lessons and reuses operators across problems as an implemented architectural fact.

Given that substrate, the next mechanical question is the \emph{transfer relation itself}: under what conditions may an episode, lesson or operator recorded in one problem fibre move into another? The tempting answer---semantic similarity---is known to be unsafe: surface vocabulary can be shared where the load-bearing structure differs, and disjoint where it agrees. The saturation results of the companion paper sharpen the demand: once the epistemic view of a research programme reaches bounded saturation, further progress must come from moving validated structure \emph{across} fibres rather than re-deriving it within one, so the transfer relation becomes the binding constraint on reuse. What is needed is a relation with an explicit scope, an explicit refusal mode, and instruments capable of falsifying claims made about it.

This paper reports what has actually been earned on that question, and no more. Three contributions survive our own audits; a fourth question is left explicitly open.

\begin{enumerate}[leftmargin=*]
\item \textbf{A specification (Section~\ref{sec:contract}).} The applicability contract: a directional structural witness binding quantity of interest (QoI), role/relation mapping, preserved invariants, explicitly non-preserved properties, target boundaries, evidence and uncertainty, evaluated by a non-compensatory three-verdict decision (\texttt{LICENSED} / \texttt{REJECTED} / \texttt{CANNOT\_CHECK}). This is a specification claim with an exact scope, not an empirical superiority claim. Section~\ref{sec:parents} places it against its acknowledged formal parent---causal transportability---and narrows the residual novelty to the four conjuncts that survive that comparison.
\item \textbf{A measured adapted method transfer (Section~\ref{sec:typed-refusal}).} The parent's strongest property is that refusal is a \emph{certificate of impossibility}. We show by construction that importing this property faithfully into the witness setting is wrong---the faithful import manufactures false impossibility certificates on every repairable, open-world and budget-limited arm of a frozen eight-arm design---and that an adapted typed refusal, which separates \texttt{MERELY\_UNLICENSED} from \texttt{CERTIFIABLY\_IMPOSSIBLE}, is correct on every runnable arm under hard gates that could have failed and did not.
\item \textbf{A methods contribution: instrument falsifiability (Section~\ref{sec:falsifiability}).} Three of our own transfer instruments were audited after passing, or while passing, their registered gates. One confirmatory study \emph{passed every frozen gate} and was then shown to be non-falsifiable; one perfect score was shown to measure serialization fidelity; one successor instrument, built to the resulting battery, was falsifiable, read its text, rejected its negative controls---and still failed its error-diversity gate, isolating its measured performance to template inversion. We preserve all three negatives verbatim and distill the battery (gold-arm distinctness, seed-failability, sampled discriminating coordinates, text-destruction probes, matched-pair attribution) into design requirements for transfer benchmarks. We regard these negatives, and the battery they validate, as the paper's most reusable result.
\item \textbf{An open question, now with its blocker measured (Sections~\ref{sec:natural-domain} and~\ref{sec:open-programme}).} Whether the contract's coordinates can be recovered from natural artifacts, and whether witness gating adds value on independently labelled natural-domain cases, is \emph{not established}. The frozen 16-item natural-domain packet is registered as power-limited for the 0.05 paired-Brier minimum detectable effect. The third-party-label route was executed once under a battery-compliant instrument on 1{,}542 externally labelled narrative pairs (ARN): the admission gate passed, the battery passed, and the outcome is an honest negative---\path{NEGATIVE__CAPABILITY_ABSENT}---which relocates the residual from label availability to a measured extraction-capability gap (Section~\ref{sec:arn-negative}). No result reported here substitutes for a capable, admissible reducer.
\end{enumerate}

Throughout, same-context analysis---including our own audit batteries---is labelled as such and is not represented as independent review. Negative results are never rewritten; successors are new versioned objects.

\subsection{Terminology at a glance}

For readers from machine learning or systems engineering, the following plain-language glosses fix the vocabulary used throughout; formal definitions appear at first technical use.

\begin{description}[leftmargin=1.4em,style=nextline]
\item[Quantity of interest (QoI)] the specific target quantity a transfer is meant to preserve; correctness is judged relative to it, not to overall similarity.
\item[Directional structural witness] the typed object that licenses moving structure from a source to a target: a role map, the invariants that must be preserved, the properties that are \emph{not} preserved, target boundary conditions, and evidence.
\item[Obligation set] the witness expanded into an executable checklist of transfer obligations, each evaluated \texttt{SATISFIED} / \texttt{VIOLATED} / \texttt{UNKNOWN}.
\item[Non-compensatory] a single violated load-bearing obligation forces rejection; strength elsewhere cannot buy it back---there is no weighted average.
\item[\texttt{LICENSED} / \texttt{REJECTED} / \texttt{CANNOT\_CHECK}] the three fail-closed verdicts: transfer allowed, transfer refused on a violation, or abstention when required evidence is missing (\texttt{UNKNOWN}~$\neq$~\texttt{VIOLATED}).
\item[\texttt{MERELY\_UNLICENSED} / \texttt{CERTIFIABLY\_IMPOSSIBLE}] the typed split of rejection: this witness failed, versus no admissible witness can succeed---the latter only under a closed-world declaration and a completed bounded exhaustion.
\item[Fibre] the retrieval index/representation scoped to a fixed (QoI, context) into which an episode is moved.
\item[Forbidden loss] a property whose loss under transfer invalidates the reuse even if the mapping otherwise fits.
\item[Demoted evidence] an evidence path deliberately denied confirmatory/independent authority (e.g.\ an AI operator standing in for an absent independent annotator).
\item[Template inversion] the failure mode in which an instrument's score measures how much of its surface its author templated, because the extractor inverts the renderer's own constructions.
\item[MDE] minimum detectable effect---the smallest effect size a frozen test is powered to detect.
\end{description}
